\section{Green's function method for Poisson's equation}

\subsection{Green's identities}

Let $u(\mathbf{r})$ and $v(\mathbf{r})$ have continuous first derivatives throughout a spatial region $T$ and on its closed boundary surface $\Sigma$, and continuous second derivatives inside $T$. The divergence theorem applied to the vector field $u\nabla v$ gives
\begin{equation}
\iint_{\Sigma} u\,\nabla v\cdot\mathrm{d}\mathbf{S}
=\iiint_{T}\nabla\cdot(u\nabla v)\,\mathrm{d}V
=\iiint_{T}u\,\Delta v\,\mathrm{d}V
+\iiint_{T}\nabla u\cdot\nabla v\,\mathrm{d}V.
\end{equation}
This is \termidx{Green's first identity}{Green's identity!first}. Interchanging $u$ and $v$ produces the companion identity
\begin{equation}
\iint_{\Sigma} v\,\nabla u\cdot\mathrm{d}\mathbf{S}
=\iiint_{T}v\,\Delta u\,\mathrm{d}V
+\iiint_{T}\nabla u\cdot\nabla v\,\mathrm{d}V.
\end{equation}
Subtracting the two relations cancels the gradient products and yields \termidx{Green's second identity}{Green's identity!second}
\begin{equation}
\label{eq:green2}
\iint_{\Sigma}\Bigl(
u\frac{\partial v}{\partial n}-v\frac{\partial u}{\partial n}
\Bigr)\,\mathrm{d}S
=\iiint_{T}\bigl(u\,\Delta v-v\,\Delta u\bigr)\,\mathrm{d}V,
\end{equation}
where $\partial/\partial n$ denotes the \emph{outward} normal derivative on $\Sigma$.

\subsection{Poisson boundary-value problem and point sources}

Consider the Poisson equation
\begin{equation}
\label{eq:poisson}
\Delta u=f(\mathbf{r}),\qquad\mathbf{r}\in T,
\end{equation}
subject to the linear boundary condition
\begin{equation}
\label{eq:bc-gen}
\Bigl[\alpha\frac{\partial u}{\partial n}+\beta u\Bigr]_{\Sigma}=\varphi(M).
\end{equation}
The three classical cases are:
\begin{itemize}
\item \textbf{Dirichlet} (first kind): $\alpha=0$, $\beta\neq0$;
\item \textbf{Neumann} (second kind): $\alpha\neq0$, $\beta=0$;
\item \textbf{Robin} (third kind): $\alpha\neq0$, $\beta\neq0$.
\end{itemize}

To represent an arbitrary source density $f$ as a continuum of point sources, introduce the Dirac delta function. Let $v(\mathbf{r},\mathbf{r}_0)$ be the field at $\mathbf{r}$ produced by a unit point source at $\mathbf{r}_0$. Then
\begin{equation}
\label{eq:v-delta}
\Delta v(\mathbf{r},\mathbf{r}_0)=\delta(\mathbf{r}-\mathbf{r}_0).
\end{equation}
Formally combining \eqref{eq:poisson} and \eqref{eq:v-delta} with Green's second identity would give
\begin{equation}
\iiint_{T}\bigl(v\Delta u-u\Delta v\bigr)\,\mathrm{d}V
=\iiint_{T}vf\,\mathrm{d}V
-\iiint_{T}u\,\delta(\mathbf{r}-\mathbf{r}_0)\,\mathrm{d}V.
\end{equation}
The difficulty is that $\Delta v$ is singular at $\mathbf{r}=\mathbf{r}_0$, so Green's identity cannot be applied on the whole of $T$.

\subsection{Excision of the singularity}

Remove a small ball $K_\varepsilon$ of radius $\varepsilon$ centered at $\mathbf{r}_0$, with spherical boundary $\Sigma_\varepsilon$. On the excised domain $T\setminus K_\varepsilon$ the identity \eqref{eq:green2} is legitimate:
\begin{equation}
\begin{aligned}
&\iiint_{T\setminus K_\varepsilon}\bigl(v\Delta u-u\Delta v\bigr)\,\mathrm{d}V\\
&\qquad=
\iint_{\Sigma}\Bigl(v\frac{\partial u}{\partial n}-u\frac{\partial v}{\partial n}\Bigr)\,\mathrm{d}S
+\iint_{\Sigma_\varepsilon}\Bigl(v\frac{\partial u}{\partial n}-u\frac{\partial v}{\partial n}\Bigr)\,\mathrm{d}S.
\end{aligned}
\end{equation}
Since $\delta(\mathbf{r}-\mathbf{r}_0)=0$ for $\mathbf{r}\neq\mathbf{r}_0$, the volume integral reduces to $\iiint_{T\setminus K_\varepsilon}vf\,\mathrm{d}V$, and therefore
\begin{equation}
\label{eq:excised}
\begin{aligned}
&\iint_{\Sigma}\Bigl(v\frac{\partial u}{\partial n}-u\frac{\partial v}{\partial n}\Bigr)\,\mathrm{d}S
+\iint_{\Sigma_\varepsilon}\Bigl(v\frac{\partial u}{\partial n}-u\frac{\partial v}{\partial n}\Bigr)\,\mathrm{d}S\\
&\qquad=
\iiint_{T\setminus K_\varepsilon}vf\,\mathrm{d}V.
\end{aligned}
\end{equation}

Near $\mathbf{r}_0$ the singular part of $v$ behaves as the Coulomb potential of charge $-\varepsilon_0$,
\begin{equation}
v\sim-\frac{1}{4\pi\,|\mathbf{r}-\mathbf{r}_0|}.
\end{equation}
As $\varepsilon\to0$ one finds
\begin{equation}
\iint_{\Sigma_\varepsilon}v\frac{\partial u}{\partial n}\,\mathrm{d}S\to0,
\qquad
\iint_{\Sigma_\varepsilon}u\frac{\partial v}{\partial n}\,\mathrm{d}S\to-u(\mathbf{r}_0).
\end{equation}
Passing to the limit in \eqref{eq:excised} yields the fundamental integral formula
\begin{equation}
\label{eq:fund-poisson}
\begin{aligned}
u(\mathbf{r}_0)
&=\iiint_{T}v(\mathbf{r},\mathbf{r}_0)f(\mathbf{r})\,\mathrm{d}V\\
&\quad-\iint_{\Sigma}\Bigl[
v(\mathbf{r},\mathbf{r}_0)\frac{\partial u}{\partial n}
-u(\mathbf{r})\frac{\partial v(\mathbf{r},\mathbf{r}_0)}{\partial n}
\Bigr]\,\mathrm{d}S.
\end{aligned}
\end{equation}

\begin{note}
Formula \eqref{eq:fund-poisson} appears to require both $u$ and $\partial u/\partial n$ on $\Sigma$. A standard boundary-value problem supplies only one linear combination of them. The remedy is to choose $v$ so that the unknown boundary data drop out. That special choice is the Green's function.
\end{note}

\subsection{Dirichlet, Robin, and Neumann Green's functions}

\paragraph{Dirichlet problem.}
Require the auxiliary field $G$ (replacing $v$) to vanish on the boundary,
\begin{equation}
G\big|_{\Sigma}=0.
\end{equation}
Then \eqref{eq:fund-poisson} becomes
\begin{equation}
\label{eq:dirichlet-G}
u(\mathbf{r}_0)
=\iiint_{T}G(\mathbf{r},\mathbf{r}_0)f(\mathbf{r})\,\mathrm{d}V
+\iint_{\Sigma}\phi(\mathbf{r})\frac{\partial G(\mathbf{r},\mathbf{r}_0)}{\partial n}\,\mathrm{d}S,
\end{equation}
where $\phi=u|_{\Sigma}$. Here $G$ is the Green's function of the first boundary-value problem.

\paragraph{Robin (third) problem.}
Impose the homogeneous condition of the same type as the given data,
\begin{equation}
\Bigl[\alpha\frac{\partial G}{\partial n}+\beta G\Bigr]_{\Sigma}=0.
\end{equation}
If $\bigl[\alpha\partial u/\partial n+\beta u\bigr]_{\Sigma}=\psi$, the solution formula reads
\begin{equation}
\label{eq:robin-G}
u(\mathbf{r}_0)
=\iiint_{T}Gf\,\mathrm{d}V
-\frac{1}{\alpha}\iint_{\Sigma}G(\mathbf{r},\mathbf{r}_0)\psi(\mathbf{r})\,\mathrm{d}S.
\end{equation}

\paragraph{Neumann problem.}
The naive requirements
\begin{equation}
\Delta G=\delta(\mathbf{r}-\mathbf{r}_0),
\qquad
\frac{\partial G}{\partial n}\Big|_{\Sigma}=0
\end{equation}
admit \emph{no} solution. Physically, a steady heat source inside an insulated region cannot reach equilibrium: the temperature rises without bound. One therefore introduces a \textbf{generalized Green's function} that includes a compensating uniform sink. In three dimensions,
\begin{equation}
\Delta G=\delta(\mathbf{r}-\mathbf{r}_0)-\frac{1}{V_T},
\qquad
\frac{\partial G}{\partial n}\Big|_{\Sigma}=0,
\end{equation}
where $V_T$ is the volume of $T$. In two dimensions the same construction uses $1/A_T$ with $A_T$ the area of the region.

\subsection{Reciprocity and observation-point form}

The Green's function is symmetric in its two arguments,
\begin{equation}
G(\mathbf{r},\mathbf{r}_0)=G(\mathbf{r}_0,\mathbf{r}).
\end{equation}
Interchanging the roles of the observation and source points, the Dirichlet formula becomes
\begin{equation}
\label{eq:dirichlet-obs}
u(\mathbf{r})
=\iiint_{T}G(\mathbf{r},\mathbf{r}_0)f(\mathbf{r}_0)\,\mathrm{d}V_0
+\iint_{\Sigma}\phi(\mathbf{r}_0)\frac{\partial G(\mathbf{r},\mathbf{r}_0)}{\partial n_0}\,\mathrm{d}S_0,
\end{equation}
and the Robin formula becomes
\begin{equation}
u(\mathbf{r})
=\iiint_{T}G(\mathbf{r},\mathbf{r}_0)f(\mathbf{r}_0)\,\mathrm{d}V_0
-\frac{1}{\alpha}\iint_{\Sigma}G(\mathbf{r},\mathbf{r}_0)\psi(\mathbf{r}_0)\,\mathrm{d}S_0.
\end{equation}

\begin{definition}[Physical content of $G$]
The volume integral is the superposition of interior sources; the surface integral encodes the influence of the boundary data. The function $G$ itself is the field of a unit point source subject to the corresponding homogeneous boundary condition.
\end{definition}

\begin{note}[Proof of reciprocity]
Fix $\mathbf{r}_1,\mathbf{r}_2\in T$. Excise small spheres about both points and apply Green's second identity to $u=G(\cdot,\mathbf{r}_1)$ and $v=G(\cdot,\mathbf{r}_2)$ on the remaining domain. The volume integrand and the contribution from the outer boundary $\Sigma$ vanish by the homogeneous boundary condition and by $\Delta G=\delta$. The limits on the two spheres give $G(\mathbf{r}_1,\mathbf{r}_2)=G(\mathbf{r}_2,\mathbf{r}_1)$.
\end{note}

For Laplace's equation ($f=0$) only the boundary integrals survive:
\begin{equation}
u(\mathbf{r})=\iint_{\Sigma}\varphi(\mathbf{r}_0)\frac{\partial G}{\partial n_0}\,\mathrm{d}S_0
\qquad\text{(Dirichlet)},
\end{equation}
\begin{equation}
u(\mathbf{r})=-\frac{1}{\alpha}\iint_{\Sigma}G\,\psi(\mathbf{r}_0)\,\mathrm{d}S_0
\qquad\text{(Robin)}.
\end{equation}
Thus Green's method also solves homogeneous equations once the appropriate $G$ is known.
